% Zdroj: PDF priklady-jaro-2023.pdf, fyzická str. 1-2. Značka: společná informatika I3. Zařazeno: I3.
\section{Objektový návrh}
\szzotazka{jaro 2023}{2}{Objektový návrh (bitmapový editor, polymorfismus, pluginy/DLL)}

\noindent\textbf{Zadání.}\par
Předpokládejte, že implementujeme aplikaci pro úpravu bitmapových obrázků. Uživatel
bude moci mít otevřených více obrázků najednou. Nový obrázek si uživatel otevře
v uživatelském rozhraní aplikace (UI) výběrem jména souboru -- aplikace dle přípony
souboru rozhodne o jeho formátu. U otevřeného obrázku si uživatel může postupně
vybírat některé z nabízených filtrů -- v UI aplikace si uživatel vždy vybere filtr
a jeho konkrétní nastavení, ten se uloží do seznamu vybraných filtrů. Nakonec bude
mít uživatel možnost všechny filtry najednou aplikovat. Každý otevřený soubor si
uživatel může kdykoliv uložit buď v původním nebo jiném podporovaném bitmapovém
formátu. Pro jednoduchost předpokládejte, že jednomu formátu obrázku odpovídá právě
jedna přípona souboru (tj. např. pro JPEG formát jen přípona \texttt{.jpg} a nikoliv
už \texttt{.jpeg}), a že všechny filtry mohou libovolně měnit obrazová data, ale
rozměry obrázku v pixelech vždy zůstávají zachované.

\medskip
Zatím máme v C\# připravenou níže uvedenou kostru hlavních tříd našeho programu, kde
třída \texttt{Image} reprezentuje načtená obrazová data (další metadata jako původní
jméno souboru atd. si budeme pamatovat jinde; poznámka: \texttt{[,]} v C\# znamená
dvojrozměrné pole), a metoda \texttt{Program.Main} představuje jednoduchý scénář
využití uvedených tříd:

\begin{lstlisting}[language={[Sharp]C}]
struct Color { public byte R, G, B; }

static class Filters {
    public static void ApplySharpness(Image image, float amount, float radius) { ... }
    public static void ApplyBrightness(Image image, float amount) { ... }
}

abstract class Image {
    public Color[,] bitmap = null;
    public abstract void Load(string fileName);
    public abstract void Save(string fileName);
}

class PngImage : Image {
    public override void Load(string fileName) { ... }
    public override void Save(string fileName) { ... }
}

class JpegImage : Image {
    public override void Load(string fileName) { ... }
    public override void Save(string fileName) { ... }
}

static class ImageLoader {
    public static Image Load(string fileName) {
        Image image;
        if (fileName.EndsWith(".png")) {
            image = new PngImage();
        } else if (fileName.EndsWith(".jpg")) {
            image = new JpegImage();
        } else {
            throw new ArgumentException();
        }
        image.Load(fileName);
        return image;
    }
}

class Program {
    public static void Main() {
        var fileName = "a.jpg";
        var image = ImageLoader.Load(fileName);
        Filters.ApplySharpness(image, 9.0f, 1.3f);
        Filters.ApplyBrightness(image, 1.7f);
        image.Save(fileName);
    }
}
\end{lstlisting}

Přepracujte celý objektový návrh tak, aby lépe umožňoval implementovat, udržovat
a dále rozšiřovat uvedenou aplikaci. Novou kostru kódu zapište v C\#, C++, nebo Javě.
Soustřeďte se na návrh typů a jejich vztahů (vše zapisujte v syntaxi zvoleného
jazyka), deklaraci metod a pouze klíčových vlastností (properties) a datových položek.
Implementaci většiny metod uvádět nemusíte -- pouze upravte implementaci
\texttt{Program.Main} tak, aby vhodně využívala váš OO návrh. Ve svém řešení uveďte
také implementaci kódu, který bude ve vašem návrhu implementovat posloupnost kroků
od jména obrazového souboru k výběru formátu souboru a načtení obrázku v tomto formátu.

\medskip
Ve svém návrhu zohledněte nejen snadnou možnost implementace všech aspektů výše
načrtnutého scénáře používání aplikace, ale také to, že do budoucna budeme chtít
přidávat velké množství různých filtrů -- implementaci nového filtru ale vždy
plánujeme doplňovat do hlavních zdrojových kódů aplikace; nepředpokládáme, že bychom
nové filtry dodávali ve formě pluginů. Tedy ve svém řešení zohledněte pouze možnost
batchové aplikace předvybraných filtrů, ale nemusíte řešit žádnou komplikovanou
infrastrukturu pro jejich obecný popis.

\medskip
Formou pluginů naopak budeme chtít přidávat podporu pro nové obrazové formáty.
Předpokládejte tedy, že každý nově podporovaný obrazový formát budeme implementovat
ve formě nového pluginu naší aplikace (který budeme za běhu načítat jako dynamicky
linkovanou knihovnu) -- chceme tedy, aby proces načítání obrázků ze souboru byl
dostatečně flexibilní a umožňoval nám za běhu dodat odkaz na nový formát souborů,
který bychom nalezli v nějakém pluginu (samotné načítání pluginů a podporu pro načítání
dynamicky linkovaných knihoven, ani přímo hledání popisu formátu v pluginu neřešte).

\medskip
\noindent\textbf{Náčrt řešení.}\par
\textit{(V~PDF bez náčrtu řešení; náčrt doplněn k~výkladu okruhu: vlastní vzorové řešení v~C++, ověřené překladem.)}

\medskip
Slabiny původní kostry: filtry jsou \texttt{static} metody (nejdou rozšiřovat bez zásahu do třídy \texttt{Filters}), výběr formátu je napevno v~\texttt{ImageLoader} (nejde dodat formát z~pluginu za běhu) a \texttt{Image} míchá obrazová data se čtením/zápisem. Přepracovaný návrh oddělí tři odpovědnosti do rozhraní a zavede registr formátů, do něhož smí plugin přidat formát za běhu.

\textbf{(1) Obrazová data jako hodnotový typ} \texttt{Bitmap} (rozměry se nemění):
\begin{lstlisting}[language=C++]
struct Color { unsigned char r, g, b; };

class Bitmap {                 // hodnotovy typ: drzi jen pixely
    int w_, h_;
    std::vector<Color> px_;    // w_*h_ pixelu, rozmery se nemeni
public:
    Bitmap(int w, int h) : w_(w), h_(h), px_(w * h) {}
    int width()  const { return w_; }
    int height() const { return h_; }
    Color& at(int x, int y)             { return px_[y * w_ + x]; }
    const Color& at(int x, int y) const { return px_[y * w_ + x]; }
};
\end{lstlisting}

\textbf{(2) Filtr a formát jako rozhraní} (entitní typy, drží se přes \texttt{unique\_ptr}):
\begin{lstlisting}[language=C++]
struct IFilter {               // jeden filtr = jedna trida
    virtual ~IFilter() = default;
    virtual void Apply(Bitmap& bmp) const = 0;
};
struct IImageFormat {          // cteni/zapis jednoho formatu
    virtual ~IImageFormat() = default;
    virtual Bitmap Load(const std::string& file) const = 0;
    virtual void   Save(const std::string& file, const Bitmap& bmp) const = 0;
};

class Sharpness : public IFilter {
    float amount_, radius_;
public:
    Sharpness(float a, float r) : amount_(a), radius_(r) {}
    void Apply(Bitmap&) const override { /* ... */ }
};
class Brightness : public IFilter {
    float amount_;
public:
    explicit Brightness(float a) : amount_(a) {}
    void Apply(Bitmap&) const override { /* ... */ }
};
class PngFormat  : public IImageFormat { /* Load, Save */ };
class JpegFormat : public IImageFormat { /* Load, Save */ };
\end{lstlisting}

\textbf{(3) Registr formátů} -- klíč je přípona, hodnota je továrna na formát; plugin zaregistruje nový formát za běhu jediným voláním:
\begin{lstlisting}[language=C++]
class FormatRegistry {
    std::map<std::string, std::function<std::unique_ptr<IImageFormat>()>> reg_;
public:
    void Register(const std::string& ext,
                  std::function<std::unique_ptr<IImageFormat>()> factory) {
        reg_[ext] = std::move(factory);     // sem smi sahnout i plugin za behu
    }
    std::unique_ptr<IImageFormat> ForFile(const std::string& file) const {
        auto it = reg_.find(file.substr(file.rfind('.')));
        if (it == reg_.end()) throw std::runtime_error("neznamy format");
        return it->second();
    }
};
\end{lstlisting}

\textbf{(4) Dokument} drží data, formát a seznam vybraných filtrů:
\begin{lstlisting}[language=C++]
class Document {
    Bitmap bmp_;
    std::unique_ptr<IImageFormat> format_;
    std::vector<std::unique_ptr<IFilter>> filters_;
public:
    Document(Bitmap b, std::unique_ptr<IImageFormat> f)
        : bmp_(std::move(b)), format_(std::move(f)) {}
    void AddFilter(std::unique_ptr<IFilter> f) { filters_.push_back(std::move(f)); }
    void ApplyAll()                  { for (auto& f : filters_) f->Apply(bmp_); }
    void Save(const std::string& file) { format_->Save(file, bmp_); }  // v puvodnim formatu
    // ulozeni v jinem formatu: format vybran podle cilove pripony (symetricke k ForFile)
    void SaveAs(const FormatRegistry& reg, const std::string& file) { reg.ForFile(file)->Save(file, bmp_); }
};
\end{lstlisting}

\textbf{(5) Scénář} (cesta od jména souboru přes výběr formátu k načtení a úpravám):
\begin{lstlisting}[language=C++]
int main() {
    FormatRegistry reg;
    reg.Register(".png", []{ return std::make_unique<PngFormat>(); });
    reg.Register(".jpg", []{ return std::make_unique<JpegFormat>(); });

    std::string file = "a.jpg";
    auto fmt   = reg.ForFile(file);        // dle pripony vyber format
    Bitmap bmp = fmt->Load(file);          // nacti obrazek v tom formatu
    Document doc(std::move(bmp), std::move(fmt));
    doc.AddFilter(std::make_unique<Sharpness>(9.0f, 1.3f));
    doc.AddFilter(std::make_unique<Brightness>(1.7f));
    doc.ApplyAll();
    doc.Save("a.jpg");           // v puvodnim formatu
    doc.SaveAs(reg, "b.png");    // nebo v jinem podporovanem formatu
}
\end{lstlisting}

Nový filtr = nová třída implementující \texttt{IFilter} (přidá se do kódu, nic dalšího se nemění). Nový formát = nová třída implementující \texttt{IImageFormat} plus jeden řádek \texttt{reg.Register}, který smí zavolat i plugin za běhu (požadovaná flexibilita načítání formátů). Uložení v~jiném formátu řeší \texttt{SaveAs}, jež vybere formát podle cílové přípony přes týž registr (\texttt{Save} ukládá v~původním formátu). \texttt{Bitmap} je hodnotový typ, \texttt{IFilter}/\texttt{IImageFormat} entitní (drženy přes \texttt{unique\_ptr}); o~úklid se stará RAII.

\medskip
\begin{csalt}[rozhraní + zásuvné moduly]
Tatáž architektura je v~C\# kratší a~bezpečnější. Filtry i~formáty jsou \texttt{interface}
(bez abstraktní třídy a~virtuálního destruktoru), o~úklid se stará GC (žádné \texttt{unique\_ptr}).
Klíčová výhoda u~\emph{pluginů}: formát z~DLL se načte reflexí, ne ručním \texttt{dlopen} /
\texttt{LoadLibrary} a~továrnami:
\begin{lstlisting}[style=cs]
interface IFilter      { void Apply(Bitmap b); }
interface IImageFormat { string Extension { get; }            // napr. ".png"
                         Bitmap Load(string path); void Save(Bitmap b, string path); }

// nacteni pluginu za behu: najdi vsechny IImageFormat v DLL a zaregistruj
foreach (var asm in Directory.GetFiles("plugins", "*.dll").Select(Assembly.LoadFrom))
    foreach (var t in asm.GetTypes()
                         .Where(t => typeof(IImageFormat).IsAssignableFrom(t) && !t.IsAbstract))
    {   var f = (IImageFormat)Activator.CreateInstance(t)!;
        registry.Register(f.Extension, f);                  // pripona -> format
    }
\end{lstlisting}
Nový filtr/formát = nová třída implementující rozhraní; plugin se přidá bez rekompilace jádra.
\end{csalt}
